% Generated from locale/tr/content/incompleteness/incompleteness-provability/fixed-point-lemma.tex; do not hand-edit.


\olfileid[tr]{inc}{inp}{fix}

\olsection{Sabit Nokta Lemması}

\begin{explain}
Sabit nokta lemması, $\Th{T}$ kuramı $\Th{Q}$'yu genişletiyorsa, her
$!B(x)$ !!{formula} için $\Th{T} \Proves !A \liff !B(\gn{!A})$ olacak
biçimde !!a{sentence}~$!A$ bulunduğunu söyler. Yalancı tümcesi durumunda
$!A$'nın, $\Th{T}$ içinde ispatlanabilir biçimde, ``$\gn{!A}$ yanlıştır''
önermesine; başka deyişle $\Gn{!A}$'nın yanlış !!{sentence}'nin G\"odel
sayısı olduğu önermesine denk olmasını isterdik. İspatın fikrini anlamak için
bunu Quine'ın $!A$ hakkındaki şu gayriresmî açıklamasıyla karşılaştırmak
yararlı olacaktır: ``{}`önüne kendi tırnaklı biçimi getirildiğinde yanlışlık
verir' ifadesi, önüne kendi tırnaklı biçimi getirildiğinde yanlışlık verir.''
Bir ifadenin alınıp önüne kendi tırnaklı biçimi getirilerek bir tümce
oluşturulması işlemine, ifadeyi \emph{köşegenleştirme}; sonucuna da ifadenin
köşegenleştirilmiş biçimi denebilir. Dolayısıyla `önüne kendi tırnaklı biçimi
getirildiğinde yanlışlık verir' ifadesinin köşegenleştirilmiş biçimi,
``{}`önüne kendi tırnaklı biçimi getirildiğinde yanlışlık verir' ifadesi,
önüne kendi tırnaklı biçimi getirildiğinde yanlışlık verir'' tümcesidir.
Quine'ın yalancı tümcesinin `yanlışlık verir' ifadesinin değil, `önüne kendi
tırnaklı biçimi getirildiğinde yanlışlık verir' ifadesinin köşegenleştirilmiş
biçimi olduğuna dikkat edin. Öyleyse yalancı tümcesini üretmek için
köşegenleştirilen özelliğin kendisi de köşegenleştirme içerir!{}

Aritmetik dilinde, bir serbest değişkenli !!a{formula}'nın tırnaklı biçimini,
onun G\"odel sayısını hesaplayıp bu sayının standart rakamını serbest
değişkenin yerine koyarak oluştururuz. $!E(x)$'in köşegenleştirilmiş biçimi,
$n=\Gn{!E(x)}$ olmak üzere $!E(\num{n})$'dir. (Bundan böyle
$\num{\Gn{!E(x)}}$ ifadesini $\gn{!E(x)}$ diye kısaltalım.) $!B(x)$ ``bir
yanlışlıktır'' anlamına geliyorsa, ``önüne kendi tırnaklı biçimi getirildiğinde
yanlışlık verir'' önermesi, ``köşegenleştirilmiş biçiminin G\"odel sayısına
uygulandığında yanlışlık verir'' olurdu. $!E(x)$'i, $n$ G\"odel sayılı
!!{formula}'nın köşegenleştirilmiş biçiminin G\"odel sayısını hesaplayan
$\fn{diag}(n)$ fonksiyonu için bir $\Obj{diag}$ simgemiz olsaydı,
$!B(\Obj{diag}(x))$ diye yazabilirdik. Quine'ın yalancı tümcesi biçimi de
bunun köşegenleştirilmiş biçimi, yani $!E(\gn{!E(x)})$ ya da
$!B(\Obj{diag}(\gn{!B(\Obj{diag}(x))}))$ olurdu. Elbette $!B(x)$ artık başka
herhangi bir özellik olabilir ve aynı kuruluş yine çalışır. Eksiklik teoremi
için $!B(x)$'i ``$x$, $\Th{T}$ içinde !!{derivable} değildir'' olarak
alacağız. O zaman $!E(x)$, ``köşegenleştirilmiş biçiminin G\"odel sayısına
uygulandığında !!a{sentence} verir; bu tümce $\Th{T}$ içinde !!{derivable}
değildir'' anlamına gelecektir.

Bunu $\Th{T}$ içinde biçimselleştirmek için $\fn{diag}$'ı
biçimselleştirmenin bir yolunu bulmalıyız. $\fn{diag}(n)$ hesaplanabilir,
hatta ilkel özyinelemeli bir fonksiyondur: $n$, $!E(x)$ formülünün G\"odel
sayısıysa $\fn{diag}(n)$, $!E(\gn{!E(x)})$'in G\"odel sayısını verir.
($\gn{!E(x)}$'in, $!E(x)$'in G\"odel sayısının standart rakamı, yani
$\num{\Gn{!E(x)}}$ olduğunu anımsayın.) $\Obj{diag}$, $\Th{T}$ içinde
$\fn{diag}$ fonksiyonunu temsil eden bir fonksiyon simgesi olsaydı, $!A$'yı
$!B(\Obj{diag}(\gn{!B(\Obj{diag}(x))}))$ formülü olarak alabilirdik. Şuna
dikkat edin:
\begin{align*}
\fn{diag}(\Gn{!B(\Obj{diag}(x))}) & = 
\Gn{!B(\Obj{diag}(\gn{!B(\Obj{diag}(x))}))} \\
& = \Gn{!A}.
\end{align*}
$\Th{T}$'nin
\[
\Obj{diag}(\gn{!B(\Obj{diag}(x))}) = \gn{!A},
\]
formülünü !!{derive} edebildiğini varsayarsak
$!B(\Obj{diag}(\gn{!B(\Obj{diag}(x))})) \liff !B(\gn{!A})$ formülünü de
!!{derive} edebilir. Fakat tanım gereği sol taraf $!A$'dır.

Elbette $\Obj{diag}$ genel olarak $\Th{T}$'nin bir fonksiyon simgesi değildir;
özellikle de $\Th{Q}$'nun simgelerinden biri değildir. Ancak $\fn{diag}$
hesaplanabilir olduğundan, $\Th{Q}$ içinde bir
$!D_{\fn{diag}}(x,y)$ formülüyle \emph{temsil edilebilir}. Bu nedenle
$!B(\Obj{diag}(x))$ yazmak yerine
$\lexists[y][(!D_{\fn{diag}}(x,y) \land !B(y))]$ yazabiliriz. Yukarıda
taslağı verilen ispat bunun dışında aynı biçimde ilerler; hatta daha
$\Th{Q}$ içinde işler.
\end{explain}

\begin{lem}
\ollabel{lem:fixed-point} $!B(x)$, tek serbest değişkeni $x$ olan herhangi
bir formül olsun. O zaman $\Th{Q} \Proves !A \liff !B(\gn{!A})$ olacak
biçimde bir $!A$ tümcesi vardır.
\end{lem}


\begin{proof}
$!B(x)$ verilsin. $!E(x)$,
$\lexists[y][(!D_{\fn{diag}}(x,y) \land !B(y))]$ formülü; $!A$ ise bunun
köşegenleştirilmiş biçimi, yani $!E(\gn{!E(x)})$ formülü olsun.

$!D_{\fn{diag}}$, $\fn{diag}$'ı temsil ettiğinden ve
$\fn{diag}(\Gn{!E(x)}) = \Gn{!A}$ olduğundan, $\Th{Q}$ şunları
!!{derive} edebilir:
\begin{align}
  & !D_{\fn{diag}}(\gn{!E(x)}, \gn{!A}) \ollabel{repdiag1} \\
  & \lforall[y][(!D_{\fn{diag}}(\gn{!E(x)},y) \lif
  \eq[y][\gn{!A}])]. \ollabel{repdiag2}
\end{align}
Şimdi $\Th{Q} \Proves !A \liff !B(\gn{!A})$ olduğunu gösterelim. Yalnızca
mantığı ve $\Th{Q}$ içinde !!{derivable} olguları kullanarak gayriresmî
biçimde akıl yürüteceğiz.

Önce $!A$'yı, yani $!E(\gn{!E(x)})$'i varsayın. $!E(x)$'in tanımına
dönersek $!E(\gn{!E(x)})$'in tam olarak
\[
\lexists[y][(!D_{\fn{diag}}(\gn{!E(x)},y) \land !B(y))].
\]
olduğunu görürüz. Böyle bir $y$ ele alın. $!D_{\fn{diag}}(\gn{!E(x)},y)$
olduğundan, \olref{repdiag2} uyarınca $y=\gn{!A}$'dır. Dolayısıyla $!B(y)$'den
$!B(\gn{!A})$ elde edilir.

Şimdi $!B(\gn{!A})$'yı varsayın. \olref{repdiag1} uyarınca
\begin{align*}
& !D_{\fn{diag}}(\gn{!E(x)}, \gn{!A}) \land !B(\gn{!A}).
\intertext{Bundan}
& \lexists[y][(!D_{\fn{diag}}(\gn{!E(x)},y) \land !B(y))].
\end{align*}
çıkar. Fakat bu tam olarak $!E(\gn{!E(x)})$, yani $!A$'dır.
\end{proof}

\begin{digress}
Bunu hesaplanabilirlik kuramındaki sabit nokta lemmasının ispatıyla
karşılaştırmalısınız. Burada bir \emph{önermeyi} kendisi cinsinden tanımlamak
isterken orada bir \emph{fonksiyonu} kendisi cinsinden tanımlamak istemiştik;
bu fark bir yana, fikir gerçekten aynıdır.
\end{digress}

\begin{prob} 
  !!^a{formula}~$!A(x)$, her $!B$ !!{sentence} için $\Th{Q}
  \Proves !B \liff !A(\gn{!B})$ ise bir \emph{doğruluk tanımı}dır. Sabit nokta
  lemmasını kullanarak hiçbir !!{formula}'nın doğruluk tanımı olamayacağını
  gösterin.
\end{prob}

